% !PRINTED 421
\weberpage{421}
%<<CONT
worin, wie wir jetzt öfter tun werden,

\weq{8}{m = -D}

gesetzt ist, so daß $m$ eine positive ganze Zahl bedeutet.

Diese Funktion, die für die Folge sehr wichtig ist, heißt die \gesperrt{Klassenfunktion}. Die Wurzeln der Gleichung

\weq{9}{H_m(u) = 0}

sind die Klasseninvarianten der Diskriminante $D = -m$, und diese Gleichung heißt darum die \gesperrt{Klassengleichung}.

Der Grad $h$ der Klassengleichung ist gleich der Zahl primitiver Klassen der quadratischen Formen der Diskriminante $D$. Wenn man nun die Quotienten (6) betrachtet, so ergibt sich aus alledem:

V. \gesperrt{Die Funktion $F_n(u, u)$ ist durch $[H_m(u)]^k$ teilbar, und der Quotient ist relativ prim zu $H_m(u)$.}

Hiernach läßt sich die Funktion $F_n(u, u)$ in Faktoren zerlegen, und es ergibt sich, wenn $C$ eine Konstante ist, die gleich den Koeffizienten der höchsten Potenz $u^N$ von $u$ ist:

\weq{10}{F_n(u, u) = C\,H_m^k(u)\,H_{m'}^{k'}(u)\,H_{m''}^{k''}(u) \ \ldots,}

wenn sich das Produkt auf alle die positiven ganzen Zahlen $m, m', m'', \ldots$ erstreckt, für die Gleichung § 114, (8)

\weq{11}{4\,n = y^2 + m\,x^2}

eigentliche Lösungen zuläßt, und $k, k', k'', \ldots$ jedesmal die Anzahl dieser Lösungen bedeutet (mit der dort angegebenen Modifikation in den beiden Ausnahmefällen $D = -4$, $D = -3$).

Es ist zunächst der Grad $N$ der Gleichung $F_n(u, u)$ zu bestimmen.

Wenn wir die Repräsentanten in (5) wie in § 69 auswählen, so wird

\weq{12}{F_n[j(\tau), j(\tau)] = \Pi\left[j(\tau) - j\left(\frac{c + \partial\,\tau}{a}\right)\right].}

Hierin bedeutete $a, \partial$ alle der Bedingung $a\,\partial = n$ genügenden Paare positiver ganzer Zahlen, % [sic] source prints singular bedeutete with plural subject.
und $c$ durchläuft ein Restsystem nach dem Modul $a$, mit Ausschluß solcher Werte, die zu dem größten gemeinschaftlichen Teiler $e$ von $a$ und $\partial$ nicht relativ prim sind, so daß die Anzahl der Werte von $c$, die zu einer Zerlegung von $n$ in die beiden Faktoren $a$ und $\partial$ gehören, gleich $a\,\varphi(e){:}e$ ist. Ist nun $N$ der Grad von $F_n(u, u)$ und $C$ der Koeffizient der höchsten Potenz von $u$, so beginnt die Entwickelung
%CONT>>
